IV.1
jatyantara-parinama prakrtya-apurat

IV.2
nimittam aprayojakam prakrtinam varana-bhedas tu tata ksetrikavat

IV.3
nirmana-cittani-asmita-matra

IV.4
pravrtti-bhede prayojakam cittam ekam anekesam

IV.5
tatra dhyana-jam anasayam

IV.6
karma-asukla-akrsnam yoginas trividham itaresam

IV.7
tatas tad-vipaka-anugunanam evabhivyaktir vasananam

IV.8
jati-desa-kala vyavahitanam apyanantaryam smrti-samskarayor eka-rupatvat

IV.9
tasam anaditvam caasiso nityatvat

IV.10
hetu-phala-asraya-alambanai samgëhitatvad esam abhave tad-abhava

IV.11
atitanagatam svarupato ‘styadhva-bhedad dharmanam

IV.12
te vyaktasuksma gunatmana

IV.13
parinamaikatvad vastu-tattvam

IV.14
vastu-samye citta-bhedat tayor vibhakta pantha

IV.15
na caika-citta-tantram vastu tad apramanakam tada kim syat

IV.16
tad-uparagapeksitvac-cittasya vastu jnaatajnaatam

IV.17
sada jnaataas citta-vëttayas tat-prabhoi purusasyaparinamitvat

IV.20
cittantara-dëasye buddhi-buddher atiprasamgai smrti-samkara ca

IV.21
citer apratisamkramayas tad-akarapattau svabuddhi-samvedanam

IV.22
drasøë-dëasyoparaktam cittam sarvartham

IV.23
tad asamkhyeya-vasanabhias citram api parartham samhatya-karitvat

IV.24
viasesa-darasina atma-bhava-bhavana-vinivëttii

IV.25
tada viveka-nimnam kaivalya-prag-bharam cittam

IV.26
tac-chidresu pratyayantarani samskarebhyai

IV.27
hanam esam klesavad uktam

IV.28
prasamkhyane ‘pyakusidasya sarvatha viveka-khyater dharma-megha samadhi

IV.29
tata klesa-karma-nivrtti

IV.30
tada sarvavarana-malapetasya jnaanasyanantyaj jnaeyam alpam

IV.31
tataikëtarthanam parinama-krama-samaptir gunanam

IV.32
ksana-pratiyogi parinamaparanta-nirgrahyai kramai

IV.33
purusa-artha-sunyana gunanam pratiprasavai kaivalyam svarupa-pratisøha va citi-asakter iti
